%=============================================================================
\section{Introduction}
\label{sec:intro}
%=============================================================================

\IEEEPARstart{A}{recurring} interaction pattern places a machine in the role of
questioner and a person in the role of oracle. A troubleshooting assistant
narrows a fault by asking about symptoms; a clinical triage tool narrows a
differential by asking about presentation; a product finder narrows a
catalogue by asking about requirements; the parlour game Akinator narrows a
space of characters by asking about traits. In each case the system holds a
belief over a discrete set of hypotheses, selects a question, observes an
answer, and revises. The engineering question is which question to ask, and
the near-universal answer is the one that maximises expected information gain.

That criterion inherits its authority from a clean theoretical setting in
which the answer is observed exactly. Human answers are not observed exactly.
A user asked whether the song they have in mind is ``energetic'' may hesitate,
may apply a personal threshold, or may simply be wrong, and the system has no
way to distinguish a considered answer from a guess. What makes this more than
a routine robustness concern is that the unreliability is not spread uniformly
over the question space. Some attributes admit near-certain answers: a
listener generally knows the performing artist and the language of the lyrics.
Others are irreducibly fuzzy, and inter-annotator agreement on affective music
descriptors is documented to be modest even among trained
annotators~\cite{hu2007mood,kim2010emotion}.

\subsection{The Problem}

The observation that motivates this work is that the two properties---how well
a question divides the belief, and how reliably it can be answered---need not
be independent, and that in at least one large real catalogue their dependence
runs in the unfavourable direction. An attribute is vague when competent
people disagree about where its boundary falls. Disagreement distributes items
on both sides of that boundary. A near-even split is exactly what information
gain rewards. The criterion can therefore acquire a structural tendency to
promote the questions that are hardest to answer, and it does so silently,
because nothing in the objective refers to answerability.

Routine preprocessing sharpens the effect. Continuous descriptors are commonly
discretised at quantiles so that they can be posed as binary questions. A
median split is a $50/50$ split by construction, which is the maximum of the
binary entropy function. The discretisation step thus assigns peak apparent
informativeness to continuous perceptual features, which are the features
users answer worst, while high-cardinality or heavily skewed identifiers
generate only lopsided, low-entropy binary questions and are deprioritised
despite being answerable with near certainty.

\subsection{What We Claim, and What We Do Not}

We are careful about the scope of the empirical claim, because our own
measurements narrow it. On a $79{,}803$-track music catalogue the bias is
present and statistically detectable: attributes annotated as reliably
answerable carry roughly half the maximum information gain of the rest. On two
further catalogues drawn from different domains---mushroom identification and
dermatological triage---it is absent. That is not a negative result to be
buried; it is a confirmation of the proposed mechanism, because neither of
those catalogues contains a continuous descriptor to discretise, and the
mechanism we blame is discretisation acting on perceptual quantities. The
honest generalisation is therefore conditional, and we state it as a scope
condition rather than as a universal: the bias should be expected wherever a
catalogue mixes quantile-binned perceptual descriptors with naturally skewed
objective metadata, and should not be expected where every attribute is
natively categorical.

The absence of the bias turns out not to settle whether the correction is
worth making, which is the most useful thing the cross-domain experiment
taught us. On dermatology, where the bias is absent, reliability-aware
selection more than doubles top-1 accuracy; on mushroom, where it is equally
absent, the correction changes nothing at all. What separates them is not the
correlation but whether the answerable attributes still carry enough
information to identify the target---and that is a property of the catalogue,
computable before any user is involved.

We are equally careful about the acquisition function. Correcting the
criterion means replacing the entropy of the induced split with the mutual
information between the target and the \emph{observed} answer, under an
explicit noise model. For a per-attribute binary symmetric channel that
quantity is $\Hb(p_q \star \varepsilon_q) - \Hb(\varepsilon_q)$, which is
textbook~\cite[Ch.~7]{cover2006} and which we claim no novelty for. What we
claim is empirical and diagnostic: that the observation model in some
attribute catalogues is heterogeneous in a specific and unfavourable way, that
this is partly an artefact of preprocessing rather than of the domain, that a
learned selection policy trained on simulated interaction inherits the same
blind spot, and that correcting it requires far less knowledge about users
than one might expect---in the end, none that a deployed system cannot recover
from its own interaction logs.

\subsection{On the Operating Point}

Two parameters govern how hard this task is, and neither is a property of any
method: how many questions the system may ask, and how many items it must
choose between. Both need stating in advance, because absolute accuracies mean
nothing without them and because a reader is entitled to know they were not
chosen to flatter.

\emph{Budget.} The catalogue has an entropy floor of $16.23$ bits, and the
capacity-corrected bound of Corollary~\ref{cor:budget} puts the minimum
sufficient budget at $20.15$ questions under the reliability annotation used
throughout. A system run at $T=20$ is therefore operating essentially
\emph{on} the information-theoretic boundary, where even perfect knowledge of
the noise barely suffices. That is a stress test, and we use it as one: the
condition sweep, and every robustness experiment in
Section~\ref{sec:robustness}, is run there. The headline comparison is run at
$T=40$, roughly twice the bound and within the range a deployed identifier
would actually spend, where reliability-aware selection reaches $26.6\%$ top-1
and $44.8\%$ top-5 against classical information gain's $16.6\%$ and $27.2\%$.
We note explicitly that this is the conservative direction to have moved in:
the advantage \emph{widens} with budget, so reporting at the floor understated
our own result rather than exaggerating it.

\emph{Catalogue size.} Eighty thousand items is the largest catalogue we have
and the hardest version of the task; it is not the version most deployed
identifiers face. Because size is a task parameter, we sweep it, and report
accuracy over a grid of sizes and budgets
(Section~\ref{sec:results-size}). At a thousand items and forty questions the
same method with the same annotation identifies the target $84.2\%$ of the
time, and $88.3\%$ if allowed sixty questions. That sweep is emphatically not
cross-validation and we do not present it as such---no model is fitted, so a
subset does not estimate performance on the whole, it simply poses an easier
question, and a number obtained on a subset must be reported with the subset
size attached. What it establishes is that the advantage we report is not an
artefact of operating near the floor: it is positive in every one of the
eighteen cells, by between $6.5$ and $32.2$ percentage points, and every cell
is significant at $\alpha = 0.01$.

Alongside both we report rank-based metrics, which are not compressed against
the floor, and the budget each method needs to reach a fixed accuracy, which is
the quantity a practitioner actually allocates.

\subsection{Limitations of Existing Approaches}

Three responses to answer noise appear in the applied literature, and none
addresses the bias described above.

The first is to replace hard candidate elimination with a probabilistic belief
update, so that a contradicted candidate is downweighted rather than
discarded. This is correct and necessary, and it is where the robustness
reported by comparative studies of interactive identification actually
originates. It does not change which questions are asked.

The second is to learn a question-selection policy from simulated
interaction, typically by regressing a realised measure of progress onto
features of the state and the query. This substitutes an approximation for a
quantity that has a closed form under the model, and in the absence of an
explicit noise term the learned policy inherits the bias from its training
signal. We do not leave this as an assertion: we implement the policy, give
it feature capacity sufficient to represent the reliability weighting, train
it both with and without noise in its rollouts, and report what it does.

The third is to discard unreliable attributes altogether. This is a reasonable
practitioner heuristic and we treat it as a baseline, but it is wasteful: an
attribute answered correctly seventy per cent of the time still transmits
information, and discarding it removes that information permanently rather
than deferring its use to the point where more reliable questions have been
exhausted.

\subsection{Research Gap}

The theory needed to correct the bias is not new. Noisy search has been
studied since R\'enyi~\cite{renyi1961} and Ulam~\cite{ulam1976}, surveyed by
Pelc~\cite{pelc2002}, and placed on a decision-theoretic footing by Jedynak et
al.~\cite{jedynak2012}. Bayesian experimental design supplies the general
principle~\cite{lindley1956,chaloner1995}, and the active-learning literature
has established both that naive greedy selection degrades badly under noisy
observations and that noise-aware greedy alternatives with approximation
guarantees exist~\cite{golovin2010noisy,golovin2011adaptive}. What is missing
is not an algorithm. It is the empirical observation that in some
attribute-indexed catalogues the observation model is strongly heterogeneous
\emph{across attributes} and adversarially aligned with the naive acquisition
function, an account of where that heterogeneity comes from, and a deployable
answer to the question of what a practitioner must actually know in order to
correct it. The noisy-search literature typically assumes a single known
crossover probability; the applied literature typically assumes none at all.

\subsection{Contributions}

\begin{enumerate}
\item \textbf{A characterisation of the bias together with its scope
condition.} We measure the relationship between per-attribute information gain
and per-attribute answerability, using a test appropriate to an annotation with
three levels rather than the rank correlation that a three-valued variable
inflates the ties of, and we isolate the contribution of the discretisation
step. We then test the
mechanism by predicting and confirming the absence of the bias in two
natively categorical catalogues (Sections~\ref{sec:results-bias}
and~\ref{sec:results-domains}).

\item \textbf{A unification that removes a standing confound.} Hard
elimination is the $\varepsilon \to 0$ limit of the Bayesian update
(Proposition~\ref{prop:limit}), so comparative studies that pair an
information-gain selector with elimination and a learned selector with soft
updating vary two factors at once. We separate them explicitly
(Section~\ref{sec:baselines}).

\item \textbf{Reliability-aware information gain.} We derive
$\RAIG(q) = \Hb(p_q \star \varepsilon_q) - \Hb(\varepsilon_q)$, show that it
re-ranks rather than rescales (Proposition~\ref{prop:reversal}), derive a
capacity-based lower bound on the query budget (Corollary~\ref{cor:budget}),
formalise the discretisation mechanism
(Proposition~\ref{prop:discretisation}), and extend the criterion to admit
abstention (Section~\ref{sec:erasure}). The change is one line in an existing
selector.

\item \textbf{A crossover vector learned from interaction logs, and a
robustness result that needs none of the assumed values.} The obvious
objection to this work is that its three reliability constants are invented.
We answer it three ways. We re-run the primary comparison over an ensemble of
worlds drawn from priors on those constants, so the conclusion is reported
across plausible settings rather than at one point; we corrupt the annotation
and measure how much of it may be wrong before the correction stops paying;
and---the substantive answer---we show the constants are unnecessary. Treating
the target item as latent and running EM over the logs of an already-deployed
incumbent recovers a per-attribute crossover vector with no annotation, no
labels and no user study, and a system shipping that vector matches the
hand-annotated one exactly (Sections~\ref{sec:reliability-em},
\ref{sec:results-ensemble}, \ref{sec:results-em}
and~\ref{sec:results-corruption}).

\item \textbf{A learned-policy comparison.} We implement, train and evaluate
the learned-selector family the literature would otherwise supply only as a
citation, and report the extent to which it does and does not recover
reliability-aware behaviour (Section~\ref{sec:results-learned}).

\item \textbf{A deployment test that requires no users.} The cross-domain
results separate two properties that are easily conflated. Whether
informativeness correlates with unanswerability decides whether classical
selection is being actively misled; whether the answerable attributes still
carry enough information to finish the job decides whether a reliability-aware
selector can do anything about it. The second is computable from a catalogue in
seconds, by counting distinct attribute vectors over the reliable subset, and
it---not the correlation---predicts which of our three catalogues benefits
(Section~\ref{sec:results-synthetic}).

\item \textbf{An open, bitwise-reproducible harness} with paired sampling,
exact McNemar testing with multiplicity control, three catalogues, and the
restricted-attribute baseline that most directly threatens the method.
\end{enumerate}
